\subsubsection{Park and Song Formalization}
\label{ssub:park}
\textit{(João, C)}\quad

The advanced allocation method is inspired by Park and Song's
prediction-based resource-allocation approach
\cite{park2019predictionbased,park2023predictiveallocation}. It uses predicted
future work to influence current assignments, but implements this idea as a
rolling-epoch heuristic rather than as the original minimum-cost maximum-flow
model.

At each allocation epoch, the method considers currently enabled tasks and
predicted future case--activity pairs. Current tasks can be assigned directly
to available and authorized resources. Selecting a predicted candidate creates
a reservation, which keeps a compatible resource idle until the task is
expected to become available. Predictions are filtered against BPMN-valid
future activities before they enter the allocation problem.

For a resource $r$ and candidate $t$, the pairwise cost is

\[
C(r,t)=
\frac{\hat{p}(r,t)}{\tau}
-\alpha_w\frac{w(t)}{\tau}
-\alpha_\pi\pi(t)
+\rho(r,t),
\]

where $\hat{p}(r,t)$ is the estimated processing time, $w(t)$ is the current
waiting time, $\pi(t)$ is the task priority, and $\tau$ is the time scale.
Short processing times reduce the cost, while waiting time and priority favor
urgent tasks.

For current tasks, $\rho(r,t)=0$. For predicted candidates, it captures configurable terms such as prediction uncertainty, expected idling time,
no-show risk, and potential future delay. In the canonical configuration, only a subset of these terms is active; the
corresponding parameters are reported in App. ~\ref{app:joao-parksong-diagnostics}.

Assignments are selected jointly within each allocation epoch using the
Hungarian algorithm. Available resources form the rows, while current tasks,
predicted candidates, and idle choices form the columns. Infeasible pairs are
excluded, and each candidate can be selected at most once.

The method is therefore a reservation-based approximation of Park and Song's
approach, not a full multi-period stochastic optimizer.